\chapter{Library-Analyse und Veränderungen}

\section{Zielsetzung und Einordnung}
Die in dieser Arbeit verwendete Python-Bibliothek \enquote{GraviTrax Connect Python Library} stellt die zentrale Softwareschicht zwischen einer externen Steueranwendung und dem GraviTrax-POWER-System dar \cite{github2023gravitraxconnectlibrary}.
Ihre praktische Relevanz ergibt sich aus der Rolle als BLE-Gateway: Einerseits muss sie eine stabile Verbindung zur Bridge aufbauen und verwalten, andererseits eingehende und ausgehende Signale in einer Form bereitstellen, die in anwendungsnahen Steuerprogrammen mit vertretbarem Aufwand nutzbar ist.

Aus wissenschaftlicher Perspektive ist die Bibliothek nicht nur ein Hilfsmittel, sondern ein Untersuchungsgegenstand.
Die Qualität der späteren Murmelbahnsteuerungen hängt unmittelbar von ihrer Robustheit im asynchronen Laufzeitverhalten, der Konsistenz ihrer API sowie der Korrektheit der Signalauswertung ab.
Vor diesem Hintergrund verfolgt das Kapitel drei Ziele: Erstens wird die Bibliothek in Aufbau und Funktion analysiert, zweitens wird ihre interne Architektur auf Schichten- und Ablaufebene beschrieben, und drittens werden die im Rahmen dieser Arbeit umgesetzten Änderungen einschließlich ihrer Motivation und ihres Nutzens begründet.

\section{Struktur und Architektur der Bibliothek}
Die Bibliothek ist modular aufgebaut und trennt Konfigurationswissen, Kommunikationslogik und Anwendungscode.
Diese Trennung ist insbesondere für wissenschaftlich orientierte Weiterentwicklung relevant, da dadurch Änderungen an Protokollparametern, Laufzeitverhalten und Schnittstellenverträgen isoliert untersucht werden können.

\subsection{Modulaufbau}
Die Kernfunktionalität liegt im Paket \texttt{gravitraxconnect} mit zwei zentralen Modulen:
\begin{itemize}
	\item \texttt{gravitrax\_bridge.py}: Implementiert die komplette BLE-Kommunikation, Signalverarbeitung und Sende-/Empfangslogik.
	\item \texttt{gravitrax\_constants.py}: Definiert Protokollkonstanten (UUIDs, Farbkanäle, Statuscodes, Steintypen, Batteriemapping).
\end{itemize}

Ergänzt wird dies durch:
\begin{itemize}
	\item Anwendungsbeispiele (Ordner \texttt{examples}),
	\item mehrere Demo-Applikationen (CLI, GUI, GPIO, Timer, Reaction Game),
	\item testbezogene Skripte (hardwareabhängige Funktionstests und neue hardwareunabhängige Unit-Tests).
\end{itemize}

\subsection{Architekturebenen}
Auf konzeptioneller Ebene lässt sich die Bibliothek in vier Schichten gliedern:
\begin{enumerate}
	\item \textbf{Konstanten- und Protokollebene:} In \texttt{gravitrax\_constants.py} sind UUIDs, Statuscodes, Farbcodes und Steintypen zentral definiert.
	\item \textbf{Transportebene (BLE):} Die Bibliothek kapselt die \texttt{bleak}-Funktionalität für Scanning, Verbindungsaufbau sowie GATT-Zugriffe.
	\item \textbf{Protokoll- und Ablaufebene:} Hier werden Bytefolgen validiert, Signale konstruiert, Prüfsummen berechnet und Benachrichtigungen semantisch aufbereitet.
	\item \textbf{Anwendungsebene:} Nutzerprogramme registrieren Callbacks und steuern die Bahn ausschließlich über die öffentliche API.
\end{enumerate}

Die Schichtung folgt damit einem klaren Entkopplungsprinzip: Low-Level-BLE-Details bleiben im Bibliothekskern verborgen, während die Anwendungsebene mit semantischen Parametern (Status, Farbe, Steintyp) arbeitet.
Für die in dieser Arbeit entwickelte Steuerlogik ist dies zentral, da dadurch die Experimente primär auf Fachlogik und nicht auf Transportdetails fokussieren.

\subsection{Kernklasse \texttt{Bridge} als Zustandscontainer}
Die Klasse \texttt{Bridge} bildet das zentrale Abstraktionsobjekt.
Sie kapselt den Zustand einer BLE-Verbindung inklusive Geräteidentität, Firmware-/Hardwaredaten, Callback-Referenzen und interner Synchronisationsmechanismen.

Relevante Entwurfsmerkmale sind:
\begin{itemize}
	\item asynchrones API-Design auf Basis von \texttt{asyncio},
	\item entkoppelte Callback-Schnittstellen für Disconnect- und Notification-Ereignisse,
	\item Sperrmechanismen (\texttt{asyncio.Lock}) für kritische Bereiche bei Notification-Auswertung und Nachrichten-ID-Vergabe,
	\item zusätzliche Utility-Funktionen für Scans, Prüfsummen und Logging.
\end{itemize}

Die folgende Initialisierung zeigt exemplarisch, dass Verbindungszustand, Callback-Registrierung und Synchronisationsobjekte in einem einzigen Objekt zusammengeführt werden.

\begin{listing}[ht]
\caption{Auszug: Zustandsinitialisierung der \texttt{Bridge}-Klasse}
\label{lst:bridge-init}
\begin{minted}{python}
class Bridge:
	def __init__(self):
		self.client = None
		self.name = None
		self.firmware = None
		self.hardware = None
		self.id = 0
		self.dc_callback = None
		self.try_reconnect = False
		self.restart_notifications = True
		self.noti_callback = None
		self.__noti_lock = asyncio.Lock()
		self.__send_lock = asyncio.Lock()
		self.__address = None
		self.__next_send_id = 0
		self.__user_disconnected = False
		self.__last_notifications = []
\end{minted}
\end{listing}

Methodisch bedeutet dies, dass die Bibliothek zustandsorientiert implementiert ist.
Der Vorteil liegt in einer klaren Objektidentität für eine Bridge-Verbindung; der Nachteil ist eine stärkere Kopplung von Lebenszyklus und Laufzeitdaten innerhalb derselben Klasse.
Für die Zielanwendung dieser Arbeit überwiegt der Vorteil, weil die Steuerung einer einzelnen aktiven Bridge der dominierende Anwendungsfall ist.

\section{Funktionale Analyse}
\subsection{Verbindungsaufbau und Lebenszyklus}
Der Verbindungsaufbau erfolgt über \texttt{connect()}, wahlweise per Gerätename oder MAC-Adresse.
Intern wird zunächst ein Device-Discovery-Schritt über den BLE-Scanner durchgeführt und anschließend ein \texttt{BleakClient} initialisiert.
Nach erfolgreicher Verbindung werden Geräteinformationen (Firmware, Hardware) abgefragt.

Die Trennung zwischen \texttt{connect()}, \texttt{is\_connected()} und \texttt{disconnect()} ist klar gehalten und erlaubt eine robuste Ablaufsteuerung in asynchronen Anwendungen.
Über \texttt{dc\_callback} und optionale Reconnect-Logik kann auf Verbindungsverluste reaktiv reagiert werden.

Der folgende Auszug zeigt den grundlegenden Ablauf des Verbindungsaufbaus.

\begin{listing}[ht]
\caption{Auszug: Verbindungsaufbau per Name oder Adresse}
\label{lst:connect-flow}
\begin{minted}{python}
async def connect(self, name_or_addr="GravitraxConnect", by_name=True, timeout=25):
	if by_name:
		device = await BleakScanner.find_device_by_filter(
			lambda d, ad: d.name and d.name.lower() == name_or_addr.lower(), timeout
		)
	else:
		device = await BleakScanner.find_device_by_address(name_or_addr, timeout)

	if not isinstance(device, BLEDevice):
		return False

	self.client = BleakClient(device, disconnected_callback=self.__disconnect_callback)
	await self.client.connect(timeout=timeout)
	self.name = device.name
	self.__address = device.address
	await self.request_bridge_info()
	return True
\end{minted}
\end{listing}

Aus Architektursicht ist hervorzuheben, dass Discovery und Verbindung nicht in separate Komponenten ausgelagert sind.
Die Bibliothek priorisiert hier ein kompaktes API-Design gegenüber strikter Trennung.
Für Forschungs- und Prototypingkontexte ist dies pragmatisch, da die Nutzung mit wenigen Aufrufen möglich bleibt.

\subsection{Senden von Steuerinformationen}
Für ausgehende Kommunikation existieren drei Ebenen:
\begin{itemize}
	\item \texttt{send\_bytes()}: Low-Level-Senden beliebiger Bytefolgen,
	\item \texttt{send\_signal()}: Erzeugung protokollkonformer GraviTrax-Signale,
	\item \texttt{send\_periodic()}: wiederholtes Senden mit zeitlichem Abstand.
\end{itemize}

Ein Sendesignal wird als 7-Byte-Struktur aufgebaut (Header, Steintyp, Status, Reserved, Message-ID, Checksumme, Farbkanal).
Die Message-ID wird sequenziell oder zufällig vergeben; zusätzlich können Mehrfachsendungen (\enquote{resends}) zur Reduktion von Paketverlusten genutzt werden.

Die Signalkonstruktion erfolgt explizit als Byte-Serialisierung.
Der entsprechende Codeausschnitt macht sichtbar, dass die Message-ID-Vergabe als kritischer Abschnitt synchronisiert wird.

\begin{listing}[ht]
\caption{Auszug: Konstruktion und Versand eines GraviTrax-Signals}
\label{lst:send-signal}
\begin{minted}{python}
async def send_signal(self, status, color_channel, random_id=False, stone=STONE_BRIDGE):
	reserved = random.randrange(0, 255)
	async with self.__send_lock:
		if random_id:
			message_id = random.randrange(0, 255)
		else:
			message_id = self.__next_send_id
		self.__next_send_id = (self.__next_send_id + 1) % 256

	checksum = (MSG_DEFAULT_HEADER + stone + status + reserved + message_id) % 256
	await self.send_bytes(bytes([
		MSG_DEFAULT_HEADER, stone, status, reserved, message_id, checksum, color_channel
	]))
\end{minted}
\end{listing}

Diese Implementierung illustriert den Entwurfsansatz der Bibliothek:
Protokollwissen bleibt intern, die Anwendung übergibt semantische Parameter.
Gleichzeitig ermöglicht die API mit \texttt{send\_bytes()} eine niedrigere Zugriffsebene für Experimente, was für explorative Arbeiten mit proprietären Protokollen vorteilhaft ist.

\subsection{Empfang und Notification-Verarbeitung}
Eingehende BLE-Notifications werden über \texttt{notification\_enable()} aktiviert und intern in \texttt{\_\_notification\_handler()} verarbeitet.
Die Implementierung extrahiert mögliche GraviTrax-Signale mittels Mustererkennung, dekodiert die Felder und übergibt sie an den nutzerdefinierten Callback.

Wesentliche Schutzmechanismen sind:
\begin{itemize}
	\item Prüfsummenvalidierung zur Erkennung fehlerhafter Pakete,
	\item Duplikaterkennung über einen Verlauf zuletzt empfangener Signale,
	\item Synchronisation über einen Lock, um konkurrierende Callback-Ausführungen konsistent zu halten.
\end{itemize}

Der nächste Ausschnitt zeigt den Kern der eingehenden Signalauswertung, insbesondere die Prüfsummenprüfung und das Weiterreichen nur valider Daten.

\begin{listing}[ht]
\caption{Auszug: Validierung eingehender Notifications}
\label{lst:notification-handler}
\begin{minted}{python}
async def __notification_handler(self, characteristic, data):
	message = " ".join("{:02x}".format(x) for x in data)
	recv_signals = [m.group() for m in re.finditer(r"13(\s[0-9a-fA-F]{2}){6}", message)]

	async with self.__noti_lock:
		for recv_signal in recv_signals:
			if recv_signal in self.__last_notifications:
				continue
			self.__last_notifications.append(recv_signal)

			signal = [
				int(recv_signal[0:2], 16), int(recv_signal[3:5], 16),
				int(recv_signal[6:8], 16), int(recv_signal[9:11], 16),
				int(recv_signal[12:14], 16), int(recv_signal[15:17], 16),
				int(recv_signal[18:20], 16),
			]

			if signal[5] != calc_checksum(signal, for_received=True):
				continue

			await self.noti_callback(self, Header=signal[0], Stone=signal[1],
									 Status=signal[2], Id=signal[4], Color=signal[6])
\end{minted}
\end{listing}

Die Bibliothek folgt damit einem klaren \enquote{parse-validate-dispatch}-Muster.
Dieses Muster ist für echtzeitnahe Steuerung zweckmäßig, da fehlerhafte Pakete früh verworfen und gültige Ereignisse mit geringer zusätzlicher Latenz an die Anwendung weitergegeben werden.

\subsection{Scan- und Hilfsfunktionen}
Die Funktion \texttt{scan\_bridges()} unterstützt die Discovery verfügbarer Bridges mit Filteroptionen, Timeout und optionalem frühzeitigem Abbruch.
Die Utility-Funktionen \texttt{calc\_checksum()}, \texttt{add\_checksum()} und Logging-Helfer ergänzen die API um wiederverwendbare Basisbausteine.

Gerade \texttt{scan\_bridges()} ist architektonisch relevant, weil diese Funktion den Übergang vom ungebundenen Systemzustand (keine Verbindung) in den adressierbaren Betriebszustand vorbereitet.
Die in dieser Arbeit vorgenommene Überarbeitung verbessert hier die deterministische Laufzeitsteuerung (genauere Timeout-Behandlung und sauberes frühes Stoppen).

\section{Identifizierte Schwachstellen vor der Überarbeitung}
Bei der Analyse der Ausgangsbasis zeigten sich mehrere praxisrelevante Punkte:
\begin{enumerate}
	\item In \texttt{connect()} lieferte ein Timeout in einem Pfad keinen expliziten booleschen Rückgabewert, obwohl die Methode als \texttt{bool} konzipiert ist.
	\item Fehlerhafte Notifications (ungültige Prüfsumme) wurden protokolliert, jedoch im selben Ablauf weiterhin an den Nutzer-Callback weitergereicht.
	\item Die manuelle Lock-Steuerung im Sendepfad war unnötig fehleranfällig (Acquire/Release ohne Kontextmanager).
	\item \texttt{scan\_bridges()} verlor durch Integer-Cast Timeout-Genauigkeit und setzte den Modus \texttt{stop\_on\_hit} nicht sauber um.
	\item Logging-Handler konnten bei wiederholtem Import mehrfach angelegt werden.
	\item Dokumentation und Metadaten waren in Teilen inkonsistent (Version, Dokumentationslink, Beispiele, Mindest-Python-Version).
\end{enumerate}

Diese Schwachstellen sind nicht unabhängig voneinander.
Insbesondere inkonsistente Rückgabewerte, unzureichende Datenvalidierung und unpräzises Scan-Verhalten verstärken sich in asynchronen Steuerprogrammen gegenseitig:
Fehler werden dort typischerweise nicht sofort als Ausnahme sichtbar, sondern manifestieren sich als intermittierende Laufzeitartefakte (fehlende Triggerreaktionen, nicht reproduzierbare Zustände, uneinheitliches Anwendungsverhalten).
Die Überarbeitung zielte daher explizit auf Determinismus und beobachtbare Konsistenz.

\section{Durchgeführte Änderungen und Begründung}
Im Rahmen dieser Studienarbeit wurden die oben genannten Punkte systematisch adressiert.

\subsection{Korrekturen in der Kernlogik}
\textbf{1) Eindeutige Rückgabesemantik in \texttt{connect()}.}
Bei einem Verbindungstimeout wird nun explizit \texttt{False} zurückgegeben.
Dies stellt sicher, dass Aufruferverhalten (z.\,B. bedingte Ablaufsteuerungen) deterministisch bleibt.

\textbf{2) Verwerfen fehlerhafter Notifications.}
Nach erkannter Prüfsummenabweichung erfolgt nun ein \texttt{continue} im Handler.
Damit werden beschädigte Pakete konsequent verworfen und gelangen nicht mehr in die Nutzerlogik.
Die Änderung verbessert die Datenintegrität der Ereignisverarbeitung.

\textbf{3) Sicherere Synchronisation im Sendepfad.}
Die manuelle Sperrenverwaltung wurde durch \texttt{async with} ersetzt.
Dadurch wird die Freigabe der Sperre auch bei Ausnahmen automatisch sichergestellt; das Risiko von Deadlocks sinkt.

\textbf{4) Robusteres Scan-Verhalten.}
\texttt{scan\_bridges()} nutzt Timeouts nun als \texttt{float} ohne Abschneiden und implementiert \texttt{stop\_on\_hit} über ein Ereignisobjekt mit \texttt{wait\_for()}.
Gleichzeitig wird der Scanner in einem \texttt{finally}-Block stets korrekt gestoppt.
Dies verbessert sowohl Präzision als auch Ressourcenbereinigung.

\begin{listing}[ht]
\caption{Auszug: Überarbeitetes Scan-Verhalten mit frühem Abbruch}
\label{lst:scan-bridges}
\begin{minted}{python}
async def scan_bridges(name="GravitraxConnect", timeout=10.0, stop_on_hit=False):
	timeout = float(timeout)
	address_list = []
	stop_event = asyncio.Event()

	def discover_callback(device, advertisement_data):
		if (device.name == name or name == "") and device.address not in address_list:
			address_list.append(device.address)
			if stop_on_hit:
				stop_event.set()

	scanner = BleakScanner(discover_callback, "")
	await scanner.start()
	try:
		if stop_on_hit:
			try:
				await asyncio.wait_for(stop_event.wait(), timeout=timeout)
			except TimeoutError:
				pass
		else:
			await asyncio.sleep(timeout)
	finally:
		await scanner.stop()
	return address_list
\end{minted}
\end{listing}

\textbf{5) Korrektur von Fehlermeldungstexten.}
In mehreren TypeError-Pfaden wurden String-Formatierungen und Meldungstexte berichtigt, damit Diagnosen im Fehlerfall eindeutig sind.

\textbf{6) Vermeidung mehrfacher Logger-Handler.}
Die Logger-Initialisierung prüft nun, ob ein benannter Stream-Handler bereits existiert.
Damit werden doppelte Logausgaben in interaktiven oder mehrfach initialisierten Laufumgebungen vermieden.

\subsection{Dokumentation und Paketmetadaten}
Zusätzlich zur Kernlogik wurden projektbegleitende Artefakte angepasst:
\begin{itemize}
	\item Versionsanhebung auf \texttt{0.3.4} in \texttt{pyproject.toml},
	\item Eintrag eines validen Dokumentationslinks,
	\item inhaltliche Korrektur der README (Mindestversion Python 3.10, präzisere Formulierungen),
	\item Ergänzung eines Abschnitts zur Testausführung,
	\item Korrektur mehrerer Codebeispiele in der Bibliotheksdokumentation (insbesondere hinsichtlich Rückgabeverhalten der asynchronen Sendefunktionen).
\end{itemize}

\subsection{Qualitätssicherung durch neue Unit-Tests}
Zur Absicherung der Änderungen wurden hardwareunabhängige Unit-Tests ergänzt.
Diese decken insbesondere Regressionen in den kritischen Stellen ab:
\begin{itemize}
	\item korrektes \texttt{False} bei Connect-Timeout,
	\item Verwerfen falscher Prüfsummen im Notification-Handler,
	\item präzise Float-Timeout-Verwendung beim Scan,
	\item funktionales frühzeitiges Stoppen bei \texttt{stop\_on\_hit=True}.
\end{itemize}

Durch den Einsatz von Mocks (u.\,a. für Scanner- und Client-Verhalten) sind diese Tests ohne physische Bridge ausführbar und verbessern damit die Reproduzierbarkeit im Entwicklungsprozess.

Ein Beispieltest zur Prüfsummenvalidierung ist nachfolgend dargestellt.

\begin{listing}[ht]
\caption{Auszug: Unit-Test zum Verwerfen fehlerhafter Notifications}
\label{lst:unit-test-checksum}
\begin{minted}{python}
async def test_notification_handler_drops_wrong_checksum(self):
	bridge = gb.Bridge()
	received = 0

	async def callback(_bridge, **_signal):
		nonlocal received
		received += 1

	bridge.noti_callback = callback
	bad_signal = bytearray([0x13, 0x01, 0x02, 0x03, 0x04, 0x00, 0x01])
	await bridge._Bridge__notification_handler(None, bad_signal)

	self.assertEqual(received, 0)
\end{minted}
\end{listing}

Auf methodischer Ebene ist bedeutsam, dass diese Tests nicht die reale BLE-Hardware ersetzen sollen, sondern eine komplementäre Prüfebene bereitstellen.
Hardwaretests bleiben für End-to-End-Validierung notwendig; Unit-Tests sichern dagegen kontrolliert einzelne Invarianten ab.

\section{Bewertung der vorgenommenen Änderungen}
Die Änderungen adressieren keine grundlegende Neugestaltung der Bibliothek, sondern gezielte Verbesserungen in Robustheit, Korrektheit und Wartbarkeit.
Insbesondere die Konsolidierung der Rückgabesemantik, die korrekte Behandlung beschädigter Notifications und das verbesserte Scan-Handling wirken sich direkt auf die Zuverlässigkeit von anwendungsseitigen Steuerprogrammen aus.

Aus methodischer Sicht ist hervorzuheben, dass die Erweiterung um automatisierte Unit-Tests die Bibliothek von einem primär hardwarezentrierten Validierungsansatz (Funktionstests mit realer Bridge) um eine softwareseitig reproduzierbare Qualitätssicherung ergänzt.
Damit wird die Weiterentwicklung der Bibliothek auch ohne permanente Hardwareverfügbarkeit erleichtert.

Architektonisch bleibt die Bibliothek bewusst schlank und anwendungsnah.
Die vorgenommenen Änderungen erhöhen die strukturelle Qualität, ohne das API-Paradigma zu brechen.
Dies ist für bestehende Nutzercodes besonders vorteilhaft, weil Kompatibilität erhalten bleibt, während das Laufzeitverhalten berechenbarer wird.

Für die in dieser Arbeit entwickelten Murmelbahnprogramme ergeben sich daraus zwei direkte Effekte:
Zum einen sinkt die Wahrscheinlichkeit schwer reproduzierbarer Kommunikationsartefakte, zum anderen steigt die Aussagekraft experimenteller Ergebnisse, weil beobachtete Effekte eher auf die Fachlogik als auf Bibliotheksinstabilitäten zurückgeführt werden können.

\section{Fazit des Kapitels}
Die GraviTrax-Connect-Python-Library stellt eine funktional geeignete Grundlage für die in dieser Arbeit realisierten Murmelbahnsteuerungen dar.
Die durchgeführten Anpassungen verbessern die Verlässlichkeit der BLE-Kommunikation, reduzieren potenzielle Fehlerquellen im asynchronen Ablauf und erhöhen die Nachvollziehbarkeit durch konsistente Dokumentation und Tests.

Die überarbeitete Bibliotheksbasis bildet damit eine tragfähige Grundlage für die in den folgenden Kapiteln dargestellten anwendungsnahen Steuerlogiken und Versuchsaufbauten.
